\chapter{Penutup}

\section{Kerangka Berpikir}
Seluruh modul ini mengikuti satu urutan. Rekaman menjadi transkrip asli, dan transkrip asli menjadi sumber bagi semua hasil lain. Setiap turunan, baik hasil rapi, notulensi, maupun dokumen resmi, harus dicocokkan kembali dengan sumbernya sebelum dibagikan. Pandora membantu mempercepat pekerjaan, tetapi tanggung jawab atas isi tetap berada pada pengguna dan pejabat yang berwenang.

\section{Kembali ke Pertanyaan Pemantik}
Pada Bab I, Anda telah menuliskan jawaban awal atas pertanyaan tentang rapat atau kegiatan yang paling sulit dicatat ulang, siapa yang paling membutuhkan catatannya, dan bagaimana Anda tahu catatan itu benar. Bandingkan jawaban itu dengan langkah-langkah yang sudah Anda pelajari, lalu catat apa yang kini terpikirkan tetapi belum muncul pada awal pelatihan.

\section{Daftar Periksa Sebelum Membagikan Hasil}
\begin{ilustrasi}
Tiga minggu setelah pelatihan, Rina diminta menyusun laporan kegiatan lapangan dan notulensi rapat kerja untuk disampaikan kepada pimpinan unit. Ia membuka Pandora dengan langkah yang sudah ia kenal. Sebelum mengirim, ia memakai daftar periksa di bawah ini. Pada kali pertama, ia menemukan satu keputusan yang tercatat sebagai usulan dan memperbaikinya. Ia tidak lagi merasa harus mendengarkan seluruh rekaman dari awal, tetapi ia tetap membuka menit-menit kunci.
\end{ilustrasi}

Periksa setiap hal berikut sebelum hasil dibagikan kepada pihak lain.

\begin{enumerate}
  \item \textbf{Sumber.} Setiap angka, nama, dan tanggal dapat ditelusuri ke transkrip asli.
  \item \textbf{Keputusan.} Hanya hal yang benar-benar diputuskan yang ditulis sebagai keputusan.
  \item \textbf{Penyebut.} Setiap angka disertai keterangan siapa yang menyebutnya.
  \item \textbf{Bagian tidak jelas.} Tidak ada tebakan. Bagian yang tidak terdengar ditandai.
  \item \textbf{Kepala dan identitas.} Kolom kepala hanya berisi informasi yang Anda ketahui.
  \item \textbf{Pengesahan.} Dokumen yang memerlukan pengesahan tidak dibagikan sebagai dokumen final sebelum disahkan.
  \item \textbf{Keterbukaan.} Bila hasil dibuat dengan bantuan AI, catat keterangan itu sesuai ketentuan instansi.
\end{enumerate}

\section{Tindak Lanjut di Unit Kerja}
Setelah pelatihan, ada empat langkah yang dapat langsung Anda kerjakan di unit kerja masing-masing:
\begin{enumerate}
  \item menentukan satu rapat atau kegiatan pertama yang akan dicatat dengan Pandora, lalu menyepakati siapa yang memeriksa hasilnya;
  \item menyusun aturan internal sederhana tentang rekaman apa saja yang boleh diunggah, sesuai Ketentuan Penggunaan Data;
  \item membuat folder kerja bersama untuk berkas Word yang sudah diperiksa, bukan untuk rekaman mentah;
  \item menyampaikan kepada pengelola masalah yang belum tercantum dalam modul ini, agar modul dapat diperbarui.
\end{enumerate}

\begin{rangkuman}
Pandora mengubah rekaman menjadi transkrip, lalu membantu menyusun hasil turunan yang rapi dan dapat dibagikan. Transkrip asli menjadi rujukan utama, dan setiap turunan perlu diperiksa terhadap sumbernya. Notulensi dan dokumen resmi memiliki aturan yang ketat, terutama untuk keputusan, angka, nama, dan kepala dokumen. Rina memulai dari rekaman yang tidak ia kenal, dan pada akhirnya ia mampu menyusun, memeriksa, dan membagikan hasilnya sendiri dengan percaya diri.
\end{rangkuman}
